\section{Limitations}
\label{sec:limitations}

\paragraph{Answer format.}
The main tables use a reasoning-style timeline prompt. On single-moment grounding it costs the uncompressed model
about 30 mIoU against the standard one-line prompt, and with it Qwen3-VL-8B stays below the best constant interval on
ActivityNet Captions and Charades-STA. PCD helps only this kind of answer. Our gains over the uncompressed model on
QVHighlights, ActivityNet and Charades-STA are therefore relative to the prompt; the full TimeLens-Bench comparison
under the one-line prompt is in progress.

\paragraph{Minimum tokens per frame.}
Equal quotas fail when a frame pair keeps only 2--4 tokens. Spending the budget on fewer frames avoids this
(\cref{sec:frames}), but the rule we use (at least about eight tokens per frame pair) was derived from the same results
and still needs held-out validation.

\paragraph{Short, densely sampled videos.}
When uniform frames are less than a second apart, spreading tokens at random beats question-driven selection.

\paragraph{Not a streaming memory.}
The memory is rebuilt per question from the frames before the question. It is causal, but the full memory is
prefilled once per question, so peak memory and time-to-first-token are those of the uncompressed model: \method
reduces the memory the decoder reads, not the cost of building it. PCD also adds a second, prompt-only decoding
pass for time-value tokens. \TODO{cost table; write-once memory re-ranked at question time.}

\paragraph{Scope of the evaluation.}
Most results use one model and 475--500-question subsets (TimeLens-8B, NExT-GQA and the TimeLens-Bench subsets
are the exceptions), and a local judge with the official prompt. Every
hyperparameter was chosen on the S-EMBER questions; the other datasets were run once with the frozen configuration,
but a held-out S-EMBER split is still to be evaluated.

\section{Conclusion}
\label{sec:conclusion}

Temporal grounding asks a compressed video memory for something multiple-choice QA does not: a complete timeline.
Pruners designed for QA remove parts of it and fall below random selection, and the decoder's video-free prior fills
the gaps, mostly by ending intervals too early. \method keeps every time anchor with an equal, question-filled quota
and writes time values against the model's own no-video prediction. With a tenth of the visual tokens it grounds
better than the uncompressed model on S-EMBER grounded QA, for a general and a grounding-tuned model; under the
standard grounding prompt, Coverage alone keeps the accuracy of the uncompressed model with a quarter of the tokens. The floors are a lesson of their own: with a reasoning-style prompt a constant interval beats
the uncompressed model on two standard benchmarks, and grounding results under compression should be read against them.
